\documentclass[11pt]{article}
\usepackage[a4paper,margin=25mm]{geometry}
\usepackage{amsmath,amssymb,amsthm,microtype}
\usepackage[hidelinks]{hyperref}
\newcommand{\E}{\mathbb E}\newcommand{\Pp}{\mathbb P}
\newtheorem{theorem}{Theorem}\newtheorem{lemma}[theorem]{Lemma}
\title{Small noise on one probability interval: fourth moments suffice}
\author{}\date{29 October 2026}
\begin{document}\maketitle

Once acquisition has confined the sign probability and the inward endpoints,
the small-noise comparison needs only fourth moments. The entire incoming
tail is summable because its count-dependent variance grows at most linearly
with the count displacement. This removes the Poisson comparison and treats
the two probability ranges together.

\begin{theorem}\label{th:small}
Let $n\ge700$ be an integer, $3/8\le p\le93/200$, $q=1-p$, and
$Y=B+U$, where $B$ is Bernoulli of parameter $p$ and $\E(U\mid B)=0$.
Write $U_j$ for the conditional noise given $B=j$. Suppose
\[
 -3/5\le U_0\le a_0,\quad -a_1\le U_1\le3/5,
 \quad a_0,a_1\ge0,\quad a_0+a_1\le1/250.
\]
Put $v=pq$, $s=\E U^2$, $\lambda=ns$, $X=(Y-p)/\sqrt{v+s}$,
$\rho=\E|X|^3$, and $\rho_B=(1-2v)/\sqrt v$.
Assume $0<\lambda\le1/100$ and $|H|\le46/25$. Choose the integer $k$
so that
\[
 t=np+H\sqrt{v+s}=k+u,\qquad -3/10\le u\le7/10,
 \qquad\delta=k-np.
\]
Such an anchor exists and is interior. If $F_B$ is the binomial CDF and
$X_1,\ldots,X_n$ are independent copies, then for every $C\ge0$,
\begin{align*}
\Delta_C={}&\sqrt n\{\Pp(\textstyle\sum X_i\le H)-\Phi(H/\sqrt n)\}
             -C\rho\\
 &-\sqrt n\{F_B(k)-\Phi(\delta/\sqrt{nv})\}+C\rho_B
 <-\lambda/100.
\end{align*}
\end{theorem}

This is a local theorem for actual iid laws. Acquisition of the compact
probability interval, support strip and raw threshold is separate. At the
sharp positive coefficient, Schulz's strict binomial theorem turns the
comparison into a negative signed excess. Zero noise is directly Bernoulli;
there is no limiting argument in the strict estimate above.

\section{An exact conditioning identity}
Let $W_j$ be a sum of $j$ independent $U_1$ variables and $n-j$ independent
$U_0$ variables, and let $\pi_j=\Pp(\operatorname{Bin}(n,p)=j)$.
Define
\[
 f(u)=\frac{\Phi((\delta+u)/\sqrt{n(v+s)})-
                 \Phi(\delta/\sqrt{n(v+s)})}{\pi_k},\qquad
 O=\Pp(W_k>u)+f(u),
\]
\[
 T_{\rm in}=\sum_{i=1}^{n-k}\Pp(W_{k+i}\le u-i).
\]
Conditioning on the actual Bernoulli count gives exactly
\begin{align*}
\Pp(\textstyle\sum X_i\le H)-F_B(k)
={}&-\sum_{i=0}^{k}\pi_{k-i}\Pp(W_{k-i}>u+i)\\
 &+\sum_{i=1}^{n-k}\pi_{k+i}\Pp(W_{k+i}\le u-i).
\end{align*}
Outgoing events are strict and incoming events closed; all threshold atoms
are retained. Discarding only nonpositive outgoing terms, we obtain
\begin{equation}\label{eq:master}
 \Delta_C\le-m\{O-RT_{\rm in}\}+\mathcal R,
 \qquad m=\sqrt n\pi_k,
\end{equation}
whenever $\pi_{k+i}/\pi_k\le R$ on the legal incoming counts. Here
\[
\mathcal R=\sqrt n\{\Phi(\delta/\sqrt{nv})-
                \Phi(\delta/\sqrt{n(v+s)})\}+C(\rho_B-\rho).
\]

\section{Finite normalization on the common interval}
Set
\[
 v_*=15/64,\quad L=1/100,\quad D=5/3,\quad b=1/20,
 \quad T=101/10000,\quad \eta_*=2/525.
\]
We have $v\ge v_*$ and $v+s<1/4$, hence $|t-np|<23/25$.
Thus $|\delta|<D$ globally and $|\delta|<1$ when $|u|\le b$.
The segment between $p$ and $k/n$ is inside $[.37,.47]$, proving
also that $1\le k\le n-1$. On this segment
\[
 z(1-z)>23/100,\quad |D'''(z\Vert p)|<5,\quad
 |\{\log(z(1-z))\}'|<6/5,\quad 1/(12z(1-z))<3/8.
\]
Finite Stirling with $0<r_j<1/(12j)$ and entropy Taylor expansion give
\begin{equation}\label{eq:mass}
 m=\frac{\phi(0)}{\sqrt v}
       \exp\{-\delta^2/(2nv)+e\},\qquad
 n|e|<Z(D_0):=3D_0/5+3/8+D_0^3/840
\end{equation}
whenever $|\delta|\le D_0\le D$. The terms are respectively the
logarithmic prefactor, factorial remainder and entropy remainder; the last
uses $n\ge700$. The elementary finite Stirling formula has its retained
proof in the manuscript's count-comparison appendix.

At a fixed count,
\[
\log f'(z)=-\tfrac12\log(1+s/v)+\frac{\delta^2}{2nv}
                  -\frac{(\delta+z)^2}{2n(v+s)}-e.
\]
Therefore, over an integration interval with $|z|\le b_0$,
\[
n|\log f'(z)|\le\frac{11}{5}\lambda+
             \frac{2D_0b_0+b_0^2}{2v_*}+Z(D_0).
\]
The coefficient $11/5$ is larger than
$1/(2v_*)+D^2/(1400v_*^2)$. Applying $e^{-x}\ge1-x$ and
$e^x\le(1-x)^{-1}$, with $(D_0,b_0)=(1,b)$ or $(D,7/10)$,
gives
\[
\begin{array}{c|cc}
 &|u|\le1/20&-3/10\le u\le7/10\\\hline
 m\text{ lower bound}&3/4&3/4\\
 f'\text{ lower bound}&99/100&49/50\\
 f'\text{ upper bound}&101/100&51/50\\
 R\text{ upper bound}&51/50&6/5.
\end{array}
\]
These bounds hold on the whole interval from $0$ to $u$. For the mass lower
bound use \eqref{eq:mass}, $\phi(0)>199/500$ and $\sqrt v\le1/2$.
For the maximum mass, finite Stirling at a mode gives
$\sqrt n\max_j\pi_j< (41/100)/\sqrt{v_*}<17/20$.
One justification of the coefficient $41/100$ is the common mode bound on
$p\in[1/3,2/3]$: at $N=n+1\ge43$, its variance product exceeds $39v/40$,
so the mass is at most
$(2/5)\sqrt{40/39}(516/515)/\sqrt{Nv}<(41/100)/\sqrt{Nv}$.
The same bound with $N=n$ suffices here. Hence globally
$R<(17/20)/(3/4)<6/5$.
In the inner interval the first exact ratio is
\[
 \frac{\pi_{k+1}}{\pi_k}
 =1+\frac{p-\delta-1}{nv+q(\delta+1)}
 <1+\frac{93/200}{700v_*}<51/50.
\]
Every subsequent ratio is at most one, because $p-\delta<2$.

The cubic correction is favorable. The conditional supports have fixed signs
relative to $p$, and expansion using their inward cubic bounds gives
\[
 (v+s)^{3/2}(\rho-\rho_B)
 \ge s\{9/8-1/250-51/64-(17/20)s\}>0.
\]
Here $(1+x)^{3/2}\le1+3x/2+3x^2/8$ controls normalization.
Differentiating the normal CDF then proves
\begin{equation}\label{eq:repair}
 \mathcal R<2\lambda/n\quad(|u|\le b),\qquad
 \mathcal R<4\lambda/n\quad\hbox{globally}.
\end{equation}
Indeed $\phi(0)/(2v_*^{3/2})<2$, and $\delta_+<1$ or $D<2$.
No positive coefficient lower bound on $C$ is needed.

\section{The entire incoming tail from fourth moments}
Write $s_j=\E U_j^2$. Since $|s_1-s_0|\le(8/3)s$, the own-count
variance $t_0=\E W_k^2$ satisfies
\[
 .99\lambda<t_0<1.01\lambda\le T.
\]
For every legal incoming count,
\[
 t_i=\E W_{k+i}^2\le\lambda\{1+\eta_*(D_0+i)\},
 \quad D_0=1\text{ in the inner interval, }D\text{ globally}.
\]
These are actual count variances of the same conditional laws.
Independence and $|U_j|\le A=3/5$ give
\[
 \E W_{k+i}^4\le A^2t_i+3t_i^2.
\]
Put $a=1+\eta_*D_0$ and $c=\eta_*$. The elementary series bounds
\[
 \sum_{i\ge1}i^{-4}<13/12,\qquad
 \sum_{i\ge1}i^{-3}<5/4,\qquad
 \sum_{i\ge1}i^{-2}<5/3
\]
follow by summing respectively the first five, three and five terms and
bounding the remainder by an integral. Define
\[
 K(a)=A^2\left(\frac{13a}{12}+\frac{5c}{4}\right)
 +3L\left(\frac{13a^2}{12}+\frac{5ac}{2}+\frac{5c^2}{3}\right).
\]
Fourth-moment Markov at the positive distance $i-u$, followed by extension
of a nonnegative sum to all positive integers, now gives
\begin{align}
 T_{\rm in}&<\frac{K(1+c)}{(1-b)^4}\lambda<\frac{53}{100}\lambda
                         &&(|u|\le b),\label{eq:inner-tail}\\
 T_{\rm in}&<\frac{43\lambda}{100(1-u)^4}
                         &&(0\le u\le7/10),\notag\\
 T_{\rm in}&<\frac{43}{100}\lambda
                         &&(u\le0).
 \label{eq:outer-tail}
\end{align}
For the last two bounds use $K(1+cD)<43/100$. If $u\ge0$, then
$i-u\ge(1-u)i$; if $u\le0$, then $i-u\ge i$.
There is no truncation and no unproved moment allocation in these bounds.

\section{One quartic minorant controls the inner phases}
Let $\alpha=13/25$ and $M=5\alpha^4$. The polynomial
\[
 h_\alpha(x)=\frac{-x^4+3\alpha^2x^2+2\alpha^3x}{M}
                      \le\mathbf1_{(0,\infty)}(x)
\]
is a strict-event minorant. Its numerator factors as
$-x(x+\alpha)^2(x-2\alpha)$, and its positive maximum divided by $M$
is $(9+6\sqrt3)/20<1$.
For $r=zt_0$, the moment bounds $|\E W_k^3|\le At_0$ and
$\E W_k^4\le A^2t_0+3t_0^2$, combined with the normal term, give
\[
 O/t_0\ge F_{\sigma,b_0}(z,t_0),
\]
\begin{align*}
F_{\sigma,b_0}(z,t)={}&\frac{3\alpha^2-A^2-3t}{M}
 +\left(\frac{2\sigma}{5\alpha}+b_0-\frac{4At}{M}\right)z\\
 &+\frac{3\alpha^2t-6t^2}{M}z^2-\frac{t^3z^4}{M}.
\end{align*}
For $u=-r$ take $(\sigma,b_0)=(1,-101/100)$; for $u=r$ take
$(-1,99/100)$. Both functions exceed $3/5$ on
$0\le z\le3/2$, $0\le t\le T$.
Here is a finite verification that explains the whole rectangle.
The functions are concave in $t$, so it suffices to check $t=0,T$.
At either endpoint replace $z^4$ by $(9/4)z^2$ in its negative term.
The resulting negative-phase quadratic decreases on $[0,3/2]$, while
the positive-phase quadratic increases there. Their minima occur at
$z=3/2$ and $z=0$, respectively; all four rational values exceed $3/5$.

For negative phases $3t_0/2\le r\le b$, Cantelli gives
\[
 O\ge\frac{r^2}{t_0+r^2}-\frac{101}{100}r>\frac35t_0.
\]
To check the strict inequality, its derivative at $r=3t_0/2$ is positive;
the derivative first increases and then decreases. Thus the minimum is at
an endpoint. At the first endpoint divide by $t_0$ and use $t_0\le T$;
at the other use $b^2/(T+b^2)-(101/100)b>(3/5)T$.
For positive phases beyond that switch, simply
$O\ge(99/100)r\ge(297/200)t_0>(3/5)t_0$.
The switch lies below $b$ since $3T/2<b$. Consequently
\begin{equation}\label{eq:own}
                  O>\frac35t_0>\frac{59}{100}\lambda
                     \qquad(|u|\le b).
\end{equation}
These arguments include $u=0$; the minorant at zero is zero and preserves
the strict outgoing event.

\section{Two outer intervals and the final comparison}
On the negative outer interval, write $u=-r$, $b\le r\le3/10$.
Cantelli and the global derivative bound give
\[
 O\ge\frac{r^2}{T+r^2}-\frac{51}{50}r>\frac7{50}\ge14\lambda.
\]
The same derivative argument reduces this to the two endpoints: the
derivative at $b$ is positive and then has at most one downward zero.
Both rational endpoint values exceed $7/50$.
Together with \eqref{eq:outer-tail} and $R\le6/5$, this gives
$O-RT_{\rm in}>(14-(6/5)(43/100))\lambda$.

On the positive outer interval $b\le u\le7/10$,
\[
 O-RT_{\rm in}\ge\frac{49}{50}u-
                  \frac{(6/5)(43/100)L}{(1-u)^4}>\frac1{25}\ge4\lambda.
\]
The displayed function is concave in $u$; both endpoint values exceed $1/25$.
Finally, in the inner interval \eqref{eq:inner-tail} and \eqref{eq:own} give
\[
 O-RT_{\rm in}>
   \left\{\frac{59}{100}-\frac{51}{50}\frac{53}{100}\right\}\lambda.
\]
Use $m>3/4$ and \eqref{eq:repair} in \eqref{eq:master}.
The smallest of the three resulting lower margins is greater than $1/100$;
in the inner case this is the single rational inequality
\[
 \frac34\left(\frac{59}{100}-\frac{51}{50}\frac{53}{100}\right)
                     -\frac2{700}>\frac1{100}.
\]
This proves Theorem~\ref{th:small} on the full noise and phase domains.

\section*{Dependencies and evidence}
The proof uses exact iid conditioning, finite Stirling, elementary Gaussian
derivatives, independent-sum moment expansions, a polynomial minorant,
Cantelli and complete reciprocal-power sums. It uses no Pinelis inequality,
Poisson comparison or parameter certificate. The companion checker verifies
the finite rational inequalities and the polynomial reductions, with
separate actual-law controls. Those controls are not the reason for coverage
of the continuous domains.

Acquisition of the stated strip and probability interval, raw threshold,
reflection/attainment and the global finite base remain separate. Joining
this theorem to the new compact-interval five-count comparison uses the
external all-$p$ binomial theorem of Schulz (2016), Theorem1, printed page1,
\href{https://d-nb.info/1197702695/34}{original thesis}. Existing wider-domain
small-noise theorems are preserved; this is a new replacement for the
acquired $n\ge700$ route, not a claim about all their original domains.
\end{document}
